\par\noindent\textbf{(1)} 谱为E\_n=ℏω(n+1/2)，故Z=Σ\_\{n=0\}∞e\textasciicircum{}\{−βℏω(n+1/2)\}=e\textasciicircum{}\{−βℏω/2\}/(1−e\textasciicircum{}\{−βℏω\})=[2sinh(βℏω/2)]\textasciicircum{}\{-1\}。必须保留零点能以求平均总能，但零点能不贡献热容。
\par\noindent\textbf{(2)} 平均能U=−∂βlnZ=(ℏω/2)coth(βℏω/2)=ℏω(1/2+n\_B)。设x=βℏω，则热容C\_V=k\_B x² e\textasciicircum{}x/(e\textasciicircum{}x−1)²，x≫1时指数趋零，x≪1时C\_V→k\_B，恢复单一经典一维振子能量k\_BT。
\par\noindent\textbf{(3)} 由∂β²lnZ求方差Var(H)=(ℏω)² e\textasciicircum{}x/(e\textasciicircum{}x−1)²=(ℏω)²n\_B(1+n\_B)。也可用正则涨落公式Var(H)=k\_BT²C\_V核对；温度为零时热态变基态，能量方差趋零。
\par\noindent\textbf{(4)} 位置与动量分别为x̂=√(ℏ/(2mω))(a+a†)、p̂=−i√(ℏmω/2)(a−a†)，热态无异常平均⟨aa⟩=0，故(Δx)²=ℏ(2n\_B+1)/(2mω)、(Δp)²=ℏmω(2n\_B+1)/2；两标准差乘积=ℏ(2n\_B+1)/2≥ℏ/2，基态饱和。
